\documentclass[10pt,a4paper]{amsart}
\usepackage{amsmath,amssymb,amsthm,graphicx}
% Retained preamble for the Stacks Project chapter morphisms.tex (Morphisms of Schemes).
% The chapter file pinned-original-source.tex carries only \input{preamble} on its first line, so no
% preamble text lives inside the chapter itself.  The Stacks Project uses one single shared preamble for
% every chapter; the snapshot stages it verbatim in sources/stacks-chow/preamble.tex.  The block below is
% copied mechanically, line for line, from that file's lines 14-18 (the Xy-pic setup) and lines 154-217
% (the native theorem declarations together with the mathematical macros).  Nothing was added, removed or
% redefined.  The wrapper class and the amsmath/amssymb/amsthm/graphicx base packages are supplied by the
% assembler, so the class/\externaldocument part of the shared preamble is not repeated here.
% For commutative diagrams we use Xy-pic
\usepackage[all]{xy}

% We use 2cell for 2-commutative diagrams.
\xyoption{2cell}
% Theorem environments.
%
\theoremstyle{plain}
\newtheorem{theorem}[subsection]{Theorem}
\newtheorem{proposition}[subsection]{Proposition}
\newtheorem{lemma}[subsection]{Lemma}

\theoremstyle{definition}
\newtheorem{definition}[subsection]{Definition}
\newtheorem{example}[subsection]{Example}
\newtheorem{exercise}[subsection]{Exercise}
\newtheorem{situation}[subsection]{Situation}

\theoremstyle{remark}
\newtheorem{remark}[subsection]{Remark}
\newtheorem{remarks}[subsection]{Remarks}

\numberwithin{equation}{subsection}

% Macros
%
\def\lim{\mathop{\mathrm{lim}}\nolimits}
\def\colim{\mathop{\mathrm{colim}}\nolimits}
\def\Spec{\mathop{\mathrm{Spec}}}
\def\Hom{\mathop{\mathrm{Hom}}\nolimits}
\def\Ext{\mathop{\mathrm{Ext}}\nolimits}
\def\SheafHom{\mathop{\mathcal{H}\!\mathit{om}}\nolimits}
\def\SheafExt{\mathop{\mathcal{E}\!\mathit{xt}}\nolimits}
\def\Sch{\mathit{Sch}}
\def\Mor{\mathop{\mathrm{Mor}}\nolimits}
\def\Ob{\mathop{\mathrm{Ob}}\nolimits}
\def\Sh{\mathop{\mathit{Sh}}\nolimits}
\def\NL{\mathop{N\!L}\nolimits}
\def\CH{\mathop{\mathrm{CH}}\nolimits}
\def\proetale{{pro\text{-}\acute{e}tale}}
\def\etale{{\acute{e}tale}}
\def\QCoh{\mathit{QCoh}}
\def\Ker{\mathop{\mathrm{Ker}}}
\def\Im{\mathop{\mathrm{Im}}}
\def\Coker{\mathop{\mathrm{Coker}}}
\def\Coim{\mathop{\mathrm{Coim}}}

% Boxtimes
%
\DeclareMathSymbol{\boxtimes}{\mathbin}{AMSa}{"02}

%
% Macros for moduli stacks/spaces
%
\def\QCohstack{\mathcal{QC}\!\mathit{oh}}
\def\Cohstack{\mathcal{C}\!\mathit{oh}}
\def\Spacesstack{\mathcal{S}\!\mathit{paces}}
\def\Quotfunctor{\mathrm{Quot}}
\def\Hilbfunctor{\mathrm{Hilb}}
\def\Curvesstack{\mathcal{C}\!\mathit{urves}}
\def\Polarizedstack{\mathcal{P}\!\mathit{olarized}}
\def\Complexesstack{\mathcal{C}\!\mathit{omplexes}}
% \Pic is the operator that assigns to X its picard group, usage \Pic(X)
% \Picardstack_{X/B} denotes the Picard stack of X over B
% \Picardfunctor_{X/B} denotes the Picard functor of X over B
\def\Pic{\mathop{\mathrm{Pic}}\nolimits}
\def\Picardstack{\mathcal{P}\!\mathit{ic}}
\def\Picardfunctor{\mathrm{Pic}}
\def\Deformationcategory{\mathcal{D}\!\mathit{ef}}

% --- appended 2026-10-12 (Chief): the project own definitions of the reference/slogan/history
% display environments, copied mechanically from the shared Stacks preamble lines 8-12
% (\usepackage{verbatim} plus the three \newenvironment lines, which the chapter preamble keeps
% verbatim from the project preamble). Nothing added, removed or redefined: these environments
% carry no mathematical content and the project itself typesets them as comments. Without these
% three lines a counted statement whose own text uses them fails with Environment slogan undefined.
\usepackage{verbatim}
\newenvironment{reference}{\comment}{\endcomment}
\newenvironment{slogan}{\comment}{\endcomment}
\newenvironment{history}{\comment}{\endcomment}
\usepackage{cleveref}
\crefname{originallemma}{original lemma}{original lemmas}
\Crefname{originallemma}{Original lemma}{Original lemmas}
\makeatletter
\newlabel{morphisms-lemma-pullback-ample-tensor-relatively-ample}{{0892}{source}{}{sourcepointer-1}{}}
\newlabel{morphisms-lemma-pullback-ample-tensor-relatively-ample@cref}{{[originallemma][0][]0892}{[1][0][]source}{}{}{}}
\makeatother
\crefname{originallemma}{original lemma}{original lemmas}
\Crefname{originallemma}{Original lemma}{Original lemmas}
\makeatletter
\newlabel{definition-relatively-ample}{{01VH}{source}{}{sourcepointer-2}{}}
\newlabel{definition-relatively-ample@cref}{{[originallemma][0][]01VH}{[1][0][]source}{}{}{}}
\makeatother
\crefname{originallemma}{original lemma}{original lemmas}
\Crefname{originallemma}{Original lemma}{Original lemmas}
\makeatletter
\newlabel{properties-definition-ample}{{01PS}{source}{}{sourcepointer-3}{}}
\newlabel{properties-definition-ample@cref}{{[originallemma][0][]01PS}{[1][0][]source}{}{}{}}
\makeatother
\crefname{originallemma}{original lemma}{original lemmas}
\Crefname{originallemma}{Original lemma}{Original lemmas}
\makeatletter
\newlabel{properties-lemma-ample-separated}{{09MP}{source}{}{sourcepointer-4}{}}
\newlabel{properties-lemma-ample-separated@cref}{{[originallemma][0][]09MP}{[1][0][]source}{}{}{}}
\makeatother
\crefname{originallemma}{original lemma}{original lemmas}
\Crefname{originallemma}{Original lemma}{Original lemmas}
\makeatletter
\newlabel{lemma-relatively-ample-separated}{{01VI}{source}{}{sourcepointer-5}{}}
\newlabel{lemma-relatively-ample-separated@cref}{{[originallemma][0][]01VI}{[1][0][]source}{}{}{}}
\makeatother
\crefname{originallemma}{original lemma}{original lemmas}
\Crefname{originallemma}{Original lemma}{Original lemmas}
\makeatletter
\newlabel{schemes-lemma-separated-permanence}{{01KU}{source}{}{sourcepointer-6}{}}
\newlabel{schemes-lemma-separated-permanence@cref}{{[originallemma][0][]01KU}{[1][0][]source}{}{}{}}
\makeatother
\crefname{originallemma}{original lemma}{original lemmas}
\Crefname{originallemma}{Original lemma}{Original lemmas}
\makeatletter
\newlabel{properties-lemma-invert-s-sections}{{01PW}{source}{}{sourcepointer-7}{}}
\newlabel{properties-lemma-invert-s-sections@cref}{{[originallemma][0][]01PW}{[1][0][]source}{}{}{}}
\makeatother
\crefname{originallemma}{original lemma}{original lemmas}
\Crefname{originallemma}{Original lemma}{Original lemmas}
\makeatletter
\newlabel{properties-proposition-characterize-ample}{{01Q3}{source}{}{sourcepointer-8}{}}
\newlabel{properties-proposition-characterize-ample@cref}{{[originallemma][0][]01Q3}{[1][0][]source}{}{}{}}
\makeatother
\crefname{originallemma}{original lemma}{original lemmas}
\Crefname{originallemma}{Original lemma}{Original lemmas}
\makeatletter
\newlabel{properties-lemma-ample-tensor-globally-generated}{{0890}{source}{}{sourcepointer-9}{}}
\newlabel{properties-lemma-ample-tensor-globally-generated@cref}{{[originallemma][0][]0890}{[1][0][]source}{}{}{}}
\makeatother
\crefname{originallemma}{original lemma}{original lemmas}
\Crefname{originallemma}{Original lemma}{Original lemmas}
\makeatletter
\newlabel{lemma-quasi-affine-O-ample}{{0891}{source}{}{sourcepointer-10}{}}
\newlabel{lemma-quasi-affine-O-ample@cref}{{[originallemma][0][]0891}{[1][0][]source}{}{}{}}
\makeatother
\crefname{originallemma}{original lemma}{original lemmas}
\Crefname{originallemma}{Original lemma}{Original lemmas}
\makeatletter
\newlabel{properties-lemma-ample-power-ample}{{01PT}{source}{}{sourcepointer-11}{}}
\newlabel{properties-lemma-ample-power-ample@cref}{{[originallemma][0][]01PT}{[1][0][]source}{}{}{}}
\makeatother
\title{Twisting a relatively ample sheaf by a pullback of an ample sheaf}
\author{The Stacks Project Authors}
\date{stacks-project:morphisms.tex:local-snapshot-20261012; GNU FDL 1.2 or later}
\begin{document}
\maketitle
\setcounter{subsection}{1}
\noindent\textit{Editorial note.} All delivered mathematical text is the author's own, copied from the pinned Stacks Project chapter file morphisms.tex; no statement or proof step was rewritten, reordered or removed.\par\medskip
\noindent\textit{Editorial note.} The chapter file pinned-original-source.tex carries only \textbackslash{}input\{preamble\} on its first line, so no preamble text lives inside the chapter. The retained file sources/stacks-morphisms/retained-preamble.tex is a mechanical, line-for-line copy of the project's single shared preamble as staged in the snapshot. Nothing was added, removed or redefined, and the excerpt serves only as preamble, never as a span coordinate.\par\medskip
\noindent\textit{Editorial note.} Both counted spans start at the \textbackslash{}begin line of their environment, so the statement and the author's proof are delivered whole; surrounding prose of the section is deliberately outside the counted spans and is carried as standard prerequisites instead.\par\medskip
\noindent\textit{Editorial note.} The citation fields (source\_id, source\_version, authors, source\_url, license.name, source\_title) are identical to those of every other card this lane has produced from this chapter, so the batch selector treats them as one source.\par\medskip
\noindent\textit{Allowed standard prerequisites (editorial).}\begin{itemize}
\item Separation and quasi-compactness consequences of ampleness and relative ampleness. Properties, Lemma \textbackslash{}ref\{properties-lemma-ample-separated\} and Lemma \textbackslash{}ref\{lemma-relatively-ample-separated\}, cited at proof lines 8845-8849. Give separatedness of the source and target.
\item The criterion producing sections with affine nonvanishing locus. Properties, Lemma \textbackslash{}ref\{properties-lemma-invert-s-sections\}, cited at proof lines 8858-8862. Produces the sections used at each point.
\end{itemize}\medskip
\ifdefined\hypertarget\hypertarget{sourcepointer-1}{}\fi
\noindent\textit{Source reference.} Original lemma 0892: focused-seeds/AGP/sources/stacks-snapshot-20261011/tags.txt, source lines 10482--10482. The lemma's own canonical tag, quoted so that readers can resolve the label inside the delivered text; it points at the counted statement itself and carries no proof content.\par\medskip
\ifdefined\hypertarget\hypertarget{sourcepointer-2}{}\fi
\noindent\textit{Source reference.} Original lemma 01VH: focused-seeds/AGP/sources/stacks-snapshot-20261011/tags.txt, source lines 2446--2446. Navigation pointer to a statement the counted proof refers to but does not contain; nothing load-bearing is delegated to it.\par\medskip
\ifdefined\hypertarget\hypertarget{sourcepointer-3}{}\fi
\noindent\textit{Source reference.} Original lemma 01PS: focused-seeds/AGP/sources/stacks-snapshot-20261011/tags.txt, source lines 2241--2241. Navigation pointer to a statement the counted proof refers to but does not contain; nothing load-bearing is delegated to it.\par\medskip
\ifdefined\hypertarget\hypertarget{sourcepointer-4}{}\fi
\noindent\textit{Source reference.} Original lemma 09MP: focused-seeds/AGP/sources/stacks-snapshot-20261011/tags.txt, source lines 12187--12187. Navigation pointer to a statement the counted proof refers to but does not contain; nothing load-bearing is delegated to it.\par\medskip
\ifdefined\hypertarget\hypertarget{sourcepointer-5}{}\fi
\noindent\textit{Source reference.} Original lemma 01VI: focused-seeds/AGP/sources/stacks-snapshot-20261011/tags.txt, source lines 2447--2447. Navigation pointer to a statement the counted proof refers to but does not contain; nothing load-bearing is delegated to it.\par\medskip
\ifdefined\hypertarget\hypertarget{sourcepointer-6}{}\fi
\noindent\textit{Source reference.} Original lemma 01KU: focused-seeds/AGP/sources/stacks-snapshot-20261011/tags.txt, source lines 2063--2063. Navigation pointer to a statement the counted proof refers to but does not contain; nothing load-bearing is delegated to it.\par\medskip
\ifdefined\hypertarget\hypertarget{sourcepointer-7}{}\fi
\noindent\textit{Source reference.} Original lemma 01PW: focused-seeds/AGP/sources/stacks-snapshot-20261011/tags.txt, source lines 2245--2245. Navigation pointer to a statement the counted proof refers to but does not contain; nothing load-bearing is delegated to it.\par\medskip
\ifdefined\hypertarget\hypertarget{sourcepointer-8}{}\fi
\noindent\textit{Source reference.} Original lemma 01Q3: focused-seeds/AGP/sources/stacks-snapshot-20261011/tags.txt, source lines 2252--2252. Navigation pointer to a statement the counted proof refers to but does not contain; nothing load-bearing is delegated to it.\par\medskip
\ifdefined\hypertarget\hypertarget{sourcepointer-9}{}\fi
\noindent\textit{Source reference.} Original lemma 0890: focused-seeds/AGP/sources/stacks-snapshot-20261011/tags.txt, source lines 10480--10480. Navigation pointer to a statement the counted proof refers to but does not contain; nothing load-bearing is delegated to it.\par\medskip
\ifdefined\hypertarget\hypertarget{sourcepointer-10}{}\fi
\noindent\textit{Source reference.} Original lemma 0891: focused-seeds/AGP/sources/stacks-snapshot-20261011/tags.txt, source lines 10481--10481. Navigation pointer to a statement the counted proof refers to but does not contain; nothing load-bearing is delegated to it.\par\medskip
\ifdefined\hypertarget\hypertarget{sourcepointer-11}{}\fi
\noindent\textit{Source reference.} Original lemma 01PT: focused-seeds/AGP/sources/stacks-snapshot-20261011/tags.txt, source lines 2242--2242. Navigation pointer to a statement the counted proof refers to but does not contain; nothing load-bearing is delegated to it.\par\medskip
\begin{lemma}
\label{lemma-pullback-ample-tensor-relatively-ample}
Let $f : X \to Y$ be a morphism of schemes, $\mathcal{M}$
an invertible $\mathcal{O}_Y$-module, and $\mathcal{L}$ an
invertible $\mathcal{O}_X$-module.
\begin{enumerate}
\item If $\mathcal{L}$ is $f$-ample and $\mathcal{M}$
is ample, then $\mathcal{L} \otimes f^*\mathcal{M}^{\otimes a}$ is ample
for $a \gg 0$.
\item If $\mathcal{M}$ is ample
and $f$ quasi-affine, then $f^*\mathcal{M}$ is ample.
\end{enumerate}
\end{lemma}

\begin{proof}
Assume $\mathcal{L}$ is $f$-ample and $\mathcal{M}$ ample.
By assumption $Y$ and $f$ are quasi-compact (see
Definition \ref{definition-relatively-ample} and
Properties, Definition \ref{properties-definition-ample}).
Hence $X$ is quasi-compact. By Properties, Lemma
\ref{properties-lemma-ample-separated} the scheme $Y$ is
separated and by Lemma \ref{lemma-relatively-ample-separated}
the morphism $f$ is separated. Hence $X$ is separated by
Schemes, Lemma \ref{schemes-lemma-separated-permanence}.
Pick $x \in X$. We can choose $m \geq 1$
and $t \in \Gamma(Y, \mathcal{M}^{\otimes m})$ such that $Y_t$
is affine and $f(x) \in Y_t$. Since $\mathcal{L}$ restricts to an
ample invertible sheaf on $f^{-1}(Y_t) = X_{f^*t}$
we can choose $n \geq 1$ and $s \in \Gamma(X_{f^*t}, \mathcal{L}^{\otimes n})$
with $x \in (X_{f^*t})_s$ with $(X_{f^*t})_s$ affine.
By Properties, Lemma \ref{properties-lemma-invert-s-sections} part (2)
whose assumptions are satisfied by the above,
there exists an integer $e \geq 1$ and a section
$s' \in \Gamma(X, \mathcal{L}^{\otimes n} \otimes f^*\mathcal{M}^{\otimes em})$
which restricts to $s(f^*t)^e$ on $X_{f^*t}$. For any $b > 0$
consider the section $s'' = s'(f^*t)^b$ of
$\mathcal{L}^{\otimes n} \otimes f^*\mathcal{M}^{\otimes (e + b)m}$.
Then $X_{s''} = (X_{f^*t})_s$ is an affine open of $X$ containing $x$.
Picking $b$ such that $n$ divides $e + b$ we see
$\mathcal{L}^{\otimes n} \otimes f^*\mathcal{M}^{\otimes (e + b)m}$
is the $n$th power of $\mathcal{L} \otimes f^*\mathcal{M}^{\otimes a}$
for some $a$ and we can get any $a$ divisible by $m$ and big enough.
Since $X$ is quasi-compact a finite number of these affine opens
cover $X$. We conclude that for some $a$ sufficiently divisible and
large enough the invertible sheaf
$\mathcal{L} \otimes f^*\mathcal{M}^{\otimes a}$ is ample on $X$.
On the other hand, we know that $\mathcal{M}^{\otimes c}$
(and hence its pullback to $X$) is globally generated for all $c \gg 0$
by Properties, Proposition \ref{properties-proposition-characterize-ample}.
Thus $\mathcal{L} \otimes f^*\mathcal{M}^{\otimes a + c}$ is ample
(Properties, Lemma \ref{properties-lemma-ample-tensor-globally-generated})
for $c \gg 0$ and (1) is proved.

\medskip\noindent
Part (2) follows from Lemma \ref{lemma-quasi-affine-O-ample},
Properties, Lemma \ref{properties-lemma-ample-power-ample}, and
part (1).
\end{proof}

\end{document}
